\newpage
\section{线性代数与离散系统分析基础知识}
\subsection{线性代数知识补充}
首先介绍\textbf{凯莱-哈密顿定理（Cayley-Hamilton Theorem）}。对于一个\(n\)阶的方阵\(\bm{A}\)，其特征方程可以表示为：
\begin{equation*}
    \det(\bm{A}-\lambda \bm{I}) = 0 = a_{0}+a_{1}\lambda+a_{2}\lambda^2+\cdots+a_{n}\lambda^n=f(\lambda)
\end{equation*}

这是一条代数方程。而凯莱-哈密顿定理指出，方阵\(\bm{A}\)同样会满足特征方程：
\begin{equation*}
    a_{0}+a_{1}\bm{A}+\bm{A}^2+\cdots+a_{n}\bm{A}^n=f(\bm{A})=0
\end{equation*}

对上式进行变换，可以得到：
\begin{equation*}
    \bm{A}^n = -\cfrac{1}{a_{n}}(a_{0}+a_{1}\bm{A}+a_{2}\bm{A}^2+\cdots+a_{n-1}\bm{A}^{n-1})
\end{equation*}

也就是说，\(\bm{A}^n\)可以表示为\(\bm{A}^0\sim \bm{A}^{n-1}\)的线性组合结果。若对上式同时左乘\(\bm{A})\)，则可以得到：
\begin{equation*}
    \bm{A}^{n+1} = \bm{A}^n+1 = -\cfrac{1}{a_{n}}(a_{0}+a_{1}\bm{A}+a_{2}\bm{A}^2+\cdots+a_{n-1}\bm{A}^{n})
\end{equation*}

可知\(\bm{A}^{n+1}\)可以表示为\(\bm{A}^\sim \bm{A}^{n}\)的线性组合结果。由于\(\bm{A}^n\)可以表示为\(\bm{A}^0\sim \bm{A}^{n-1}\)的线性组合，
因此\(\bm{A}^{n+1}\)也可以表示为\(\bm{A}^0\sim \bm{A}^{n-1}\)的线性组合。根据数学归纳法可知：

\textbf{对\(n\)阶方阵\(\bm{A}\)，\(\forall k\ge n\)，\(\bm{A}^k\)可以表示为\(\bm{A}^0\sim \bm{A}^{n-1}\)的线性组合。}

其次介绍矩阵的\textbf{奇异值分解（Singular Value Decomposition,SVD）}，这是后续系统辨识最核心的知识。

对于\(n\)阶方阵\(\bm{A}\)，可以求取其秩为\(\textrm{rank} (\bm{A})\)，秩的值表示了\(\bm{A}\)矩阵具有多少个非零的特征值。
但对于非方阵\(\bm{A})\)必然是不满秩不能进行特征值分解的，但可以进行奇异值分解。因此奇异值分解可以视为特征值分解在任意维矩阵上的推广。

根据方阵的特征值分解公式，方阵\(\bm{A}\)可以表示为：
\begin{equation*}
    \bm{A} = \bm{P}\bm{\Lambda}\bm{P}^{-1}
\end{equation*}

其中\(\bm{P}\)为特征向量矩阵，\(\bm{\Lambda}\)为特征值组成的对角矩阵。奇异值分解说明任意矩阵\(\bm{A}\)都可以表示为：
\begin{equation}
    \bm{A} = \bm{U}\bm{\Sigma}\bm{V}^T
    \label{eq:奇异值分解}
\end{equation}

其中\(\bm{U}\)为左奇异矩阵，\(\bm{\Sigma}\)为奇异值组成的对角矩阵，\(\bm{V}\)为右奇异矩阵。需要注意的是，奇异值是非负的，
因为他表示矩阵在各个维度上的能量分布。

最后介绍\textbf{托普利茨矩阵（Toeplitz Matrix）}，这在求解离散卷积方程中非常有用。托普利茨矩阵定义为
\textbf{从左上到右下，每一个对称对角线上具有完全相同元素的矩阵}。比如一个四阶托普利茨矩阵为
\(
\bm{\mathcal{T}} = 
\begin{bmatrix}
    4 & 1 & 0 & 2 \\
    3 & 4 & 1 & 0 \\
    2 & 3 & 4 & 1 \\
    1 & 2 & 3 & 4 \\
\end{bmatrix}
\)。

特殊地，如果一个托普利茨矩阵从左下到右上，每一个对称对角线上也具有相同的元素，则称为对称托普利茨矩阵。比如
一个四阶的对称托普利茨矩阵为
\(
\bm{\mathcal{T}} = 
\begin{bmatrix}
    4 & 3 & 2 & 1 \\
    3 & 4 & 3 & 2 \\
    2 & 3 & 4 & 3 \\
    1 & 2 & 3 & 4 \\
\end{bmatrix}
\)。

\subsection{离散系统知识补充}
离散系统的状态空间方程与连续系统类似，可以表示为：
\begin{equation}
    \begin{cases}
        x_{k+1} = \bm{A} x_{k} + \bm{B} u_{k} \\
        y_{k} = \bm{C} x_{k} + \bm{D} u_{k}
    \end{cases}
    \label{eq:离散状态空间方程}
\end{equation}

利用状态空间方程\ref{eq:离散状态空间方程}，可以得到系统的传递函数：
\begin{equation}
    G(z)=\bm{C}(zI-\bm{A})^{-1}\bm{B}+\bm{D}
    \label{eq:离散传递函数}
\end{equation}

其中\(z\)为\(\mathcal{Z}\)变换的变量，与拉普拉斯变换的\(s\)变量对应，两者满足\(z=e^{Ts}\)，其中\(T\)表示离散系统的采样时间。

一个系统的传递函数，就是其脉冲响应的拉普拉斯变换。同理，在离散系统中，其传递函数是脉冲响应序列的\(\mathcal{Z}\)变换。
假设系统的脉冲响应序列为\(g_{k}\)），则其传递函数可以表示为：

\begin{equation*}
    G(z)=\sum_{k=0}^{\infty}g_{k}z^{-k}
\end{equation*}

而根据\ref{eq:离散传递函数}，可以使用类似于下式的无穷级数展开：
\begin{equation*}
    \cfrac{1}{1-x}=1+\sum_{k=1}^{\infty}x^k
\end{equation*}

对于\((z\bm{I}-\bm{A})^{-1}\)，可以表示为\(z^{-1}(\bm{I}-z^{-1}\bm{A})^{-1}\)。后者可以无穷级数展开得到：
\begin{equation*}
    (\bm{I}-z^{-1}\bm{A})^{-1} =\bm{I}+z^{-1}\bm{A}+z^{-2}\bm{A}^2+\cdots=\sum_{k=0}^{\infty}z^{-k}\bm{A}^k
\end{equation*}

带入传递函数可得：
\begin{equation*}
    G(z)=\bm{C}(z\bm{I}-\bm{A})^{-1}\bm{B}+\bm{D}=
    \bm{C}z^{-1}\sum_{k=0}^{\infty}z^{-k}\bm{A}^k\bm{B}+\bm{D}=
    \bm{D}+\sum_{k=0}^{\infty}\bm{C}\bm{A}^k\bm{B}z^{-(k+1)}=\sum_{k=0}^{\infty}g_{k}z^{-k}
\end{equation*}

因此脉冲响应序列满足：
\begin{equation}
    g_{k} = 
    \begin{cases}
        \bm{D} & k=0 \\
        \bm{C}\bm{A}^{k-1}\bm{B} & k\ge 1
    \end{cases}
    \label{eq:脉冲响应序列}
\end{equation}

这也称为离散系统的\textbf{马尔科夫系数（Markov Coefficient）}。而根据凯莱-哈密顿定理，对于\(n\)阶方阵\(\bm{A}\)（同时\(n\)也是系统阶次），
\(\forall k\ge n\)，\(\bm{A}^k\)可以表示为\(\bm{A}^0\sim \bm{A}^{n-1}\)的线性组合。因此可以得到离散系统一个重要的性质，即
\textbf{\(n\)阶离散系统的脉冲响应序列只有前\(n\)个值是独立的，后续的值都可以被前\(n\)个值线性表示}。

为了将这一个性质直观表示，进行进一步分析，可以构造离散系统的\textbf{汉克尔矩阵（Hankel Matrix）}。离散系统的汉克尔矩阵定义为：
\begin{equation}
    \bm{\mathcal{H}} =
    \begin{bmatrix}
        g_{1} & g_{2} & \cdots & g_{n} \\
        g_{2} & g_{3} & \cdots & g_{n+1} \\
        \vdots & \vdots & \ddots & \vdots \\
        g_{n} & g_{n+1} & \cdots & g_{2n-1}
    \end{bmatrix}
    \label{eq:汉克尔矩阵}
\end{equation}

根据离散系统的特性，可知\(n\)阶离散系统的汉克尔矩阵是满秩的。如果对汉克尔矩阵进行奇异值分解，也应该得到\(n\)个非零的奇异值。
